\documentclass[a4paper,12pt]{article}
\addtolength{\oddsidemargin}{-1.cm}
\addtolength{\textwidth}{2cm}
\addtolength{\topmargin}{-2cm}
\addtolength{\textheight}{3.5cm}
\newcommand{\HRule}{\rule{\linewidth}{0.5mm}}
\newcommand\tab[1][1cm]{\hspace*{#1}}
\makeindex
%%%%%%%%%%%%%%%%%%%%%%%%%%%%%%%%%%%%%%%%%%%%%%%%%%%%%%%%%%%%%%%%%%%%
\usepackage{longtable}
\usepackage{graphicx}
\usepackage{makeidx}
\usepackage{hyperref}
\usepackage{verbatim}

\hypersetup{
    colorlinks=true,
    linkcolor=black,
    filecolor=magenta,      
    urlcolor=cyan,
}
\begin{document}


% define the title
\author{Whoosh Division}
\title{OcuViz - EpiUse Labs}
\setlength{\parskip}{6pt}

% generates the title
\begin{titlepage}
\begin{center}
% Upper part of the page       
\includegraphics[scale=1]{../Spec/images/up-logo.jpg}\\[0.4cm]    
\textsc{\LARGE Department of Computer Science}\\[1.5cm]
\textsc{\Large Whoosh Division}\\[0.5cm]
% Title
\HRule \\[0.4cm]
{ \huge \bfseries OcuViz - EpiUse Labs}\\[0.4cm]
{ \huge Unit Testing Plan \& Report}\\[0.4cm]
{ Version 1.0 }\\[0.4cm]
\HRule \\[0.4cm]
\begin{minipage}{0.4\textwidth}
\begin{flushleft} \large
Vukile {Langa}
\end{flushleft}
\end{minipage}
\begin{minipage}{0.4\textwidth}
\begin{flushright} \large
\emph{} \\
u14035449 
\end{flushright}
\end{minipage}
\begin{minipage}{0.4\textwidth}
\begin{flushleft} \large
Wynand Hugo Meiring
\end{flushleft}
\end{minipage}
\begin{minipage}{0.4\textwidth}
\begin{flushright} \large
\emph{} \\
u13230795  
\end{flushright}
\end{minipage}
\begin{minipage}{0.4\textwidth}
\begin{flushleft} \large
Nontokozo Hlastwayo
\end{flushleft}
\end{minipage}
\begin{minipage}{0.4\textwidth}
\begin{flushright} \large
\emph{} \\
u14414555
\end{flushright}
\end{minipage}
\begin{minipage}{0.4\textwidth}
\begin{flushleft} \large
Gerome Schutte
\end{flushleft}
\end{minipage}
\begin{minipage}{0.4\textwidth}
\begin{flushright} \large
\emph{} \\
u12031519
\end{flushright}
\end{minipage}
\vfill

{\large \today}
\end{center}
\end{titlepage}
\footnotesize
\normalsize

%
%
%
%
%
%
%	Table of Contents
%

\tableofcontents
\newpage
\setlength{\parindent}{0em}

\section{Introduction}
\subsection{Purpose}
	OcuViz is a platform for creating rich visual representations of otherwise unintuitive data. It addresses the time consuming task of working with 3D scenes and makes it simple, by providing a format for specifying scenes that is concise, optimised for integration with pluggable data, and is targeted at being easily usable by anyone with even the slightest development experience.\\\\
	To achieve these goals, the Whoosh Division follows a test-driven development approach for OcuViz. Not only is test write-up an integral part of understanding system objects, but this allows us to ensure that all the parts of the complex system work as intended, notifying us of any breaking changes made while also identifying problematic system components early on in development, enabling us to issue targeted fixes, and giving us confidence in the soundness of system functionality at each release.\\\\
	This document combines our unit test plan and report into a single coherent artifact.
\subsection{Scope}
\subsection{Test Environment}
\subsection{Assumptions and Dependencies}

\newpage

\begin{center}
	\huge \bfseries Unit Test Plan \\[2cm]
\end{center}

\section{Test Items}

\section{Functional Features to be Tested}

\section{Test Cases}

\section{Item Pass/Fail Criteria}

\section{Test Deliverables}

\newpage

\begin{center}
	\huge \bfseries Unit Test Report \\[2cm]
\end{center}

\section{Detailed Test Results}
\subsection{Overview of Test Results}
\subsection{Functional Requirements Test Results}

\section{Other}

\section{Conclusions and Recommendations}

%
% End of Document
%
\end{document}
